\section{账本事件的区域史深描：southern-song}

\label{sec:ledger-depth-southern-song}

本组把账本事件置于战争、制度、区域社会与证据类型的交叉处。年表提供政治节点，地方材料与物质遗存则限制我们对征发、身份、信仰和迁徙所能作出的断言。

\subsection{战争、行政与区域资源}

账本条目\texttt{SS-1127-EVENT-111}所记1127年的“赵构即位为高宗，南宋朝廷开始重建”，属于南宋的历史语境。条目所示前因是前期边防、财政和统治联盟的压力构成直接背景，具体条件须按本条材料分辨，其可见后果为改变后续资源配置与政治选择；实际执行范围和社会影响仍须分项核验。这一变化若进入区域生活，必须通过道路、文书、军队、市场、家庭和地方中介才会发生；政权易手或法令发布不等于所有城镇、乡村与边地在同一时间获得同样经验。

本条所列来源为《宋史》卷24—47〈本纪〉；《建炎以来系年要录》本年条；《宋会要辑稿》相关门类。它们能够支持特定的年序、制度语言、政治行动或局部遗存，却不能无条件化为社会总量。账本特别提醒“编年主干可靠；细节、规模与动机须与对方材料、地方文书或考古材料互校”。因此，军报的兵数、条约的名分、官府的族群分类、判牍中的道德判断以及后出的叙述都须回到产生它们的场景，而不可直接代表全域动机或人口。

将事件与地方层次连接时，应同时考察资源怎样被征集、信息怎样传递、迁徙怎样改变户籍和劳动、宗教与亲属网络怎样提供或限制协作。现有证据往往只照亮少数地点；保留这种不均衡，才能避免把宋辽夏金元的多政权历史改写成单一中心的必然过程。

账本条目\texttt{SS-1127-REGNAL-YEAR}所记1127年的“南宋以宋高宗在位第1年作为诏令、赋役与外交文书的纪年框架”，属于南宋的历史语境。条目所示前因是新岁颁历、年号和纪年是中枢与地方行政沟通的共同前提，其可见后果为为同年政令、战事和地方材料提供日期锚点，不能单独推知具体政策执行。这一变化若进入区域生活，必须通过道路、文书、军队、市场、家庭和地方中介才会发生；政权易手或法令发布不等于所有城镇、乡村与边地在同一时间获得同样经验。

本条所列来源为《宋史》卷24—47〈本纪〉；《建炎以来系年要录》本年条；《宋会要辑稿》相关门类。它们能够支持特定的年序、制度语言、政治行动或局部遗存，却不能无条件化为社会总量。账本特别提醒“纪年框架可靠；该记录仅确证政权纪年连续，年内具体事项须由本年条另证”。因此，军报的兵数、条约的名分、官府的族群分类、判牍中的道德判断以及后出的叙述都须回到产生它们的场景，而不可直接代表全域动机或人口。

将事件与地方层次连接时，应同时考察资源怎样被征集、信息怎样传递、迁徙怎样改变户籍和劳动、宗教与亲属网络怎样提供或限制协作。现有证据往往只照亮少数地点；保留这种不均衡，才能避免把宋辽夏金元的多政权历史改写成单一中心的必然过程。

账本条目\texttt{SS-1128-EVENT-112}所记1128年的“高宗朝廷驻跸扬州等地，组织勤王与财政”，属于南宋的历史语境。条目所示前因是前期边防、财政和统治联盟的压力构成直接背景，具体条件须按本条材料分辨，其可见后果为改变后续资源配置与政治选择；实际执行范围和社会影响仍须分项核验。这一变化若进入区域生活，必须通过道路、文书、军队、市场、家庭和地方中介才会发生；政权易手或法令发布不等于所有城镇、乡村与边地在同一时间获得同样经验。

本条所列来源为《宋史》卷24—47〈本纪〉；《建炎以来系年要录》本年条；《宋会要辑稿》相关门类。它们能够支持特定的年序、制度语言、政治行动或局部遗存，却不能无条件化为社会总量。账本特别提醒“编年主干可靠；细节、规模与动机须与对方材料、地方文书或考古材料互校”。因此，军报的兵数、条约的名分、官府的族群分类、判牍中的道德判断以及后出的叙述都须回到产生它们的场景，而不可直接代表全域动机或人口。

将事件与地方层次连接时，应同时考察资源怎样被征集、信息怎样传递、迁徙怎样改变户籍和劳动、宗教与亲属网络怎样提供或限制协作。现有证据往往只照亮少数地点；保留这种不均衡，才能避免把宋辽夏金元的多政权历史改写成单一中心的必然过程。

账本条目\texttt{SS-1128-REGNAL-YEAR}所记1128年的“南宋以宋高宗在位第2年作为诏令、赋役与外交文书的纪年框架”，属于南宋的历史语境。条目所示前因是新岁颁历、年号和纪年是中枢与地方行政沟通的共同前提，其可见后果为为同年政令、战事和地方材料提供日期锚点，不能单独推知具体政策执行。这一变化若进入区域生活，必须通过道路、文书、军队、市场、家庭和地方中介才会发生；政权易手或法令发布不等于所有城镇、乡村与边地在同一时间获得同样经验。

本条所列来源为《宋史》卷24—47〈本纪〉；《建炎以来系年要录》本年条；《宋会要辑稿》相关门类。它们能够支持特定的年序、制度语言、政治行动或局部遗存，却不能无条件化为社会总量。账本特别提醒“纪年框架可靠；该记录仅确证政权纪年连续，年内具体事项须由本年条另证”。因此，军报的兵数、条约的名分、官府的族群分类、判牍中的道德判断以及后出的叙述都须回到产生它们的场景，而不可直接代表全域动机或人口。

将事件与地方层次连接时，应同时考察资源怎样被征集、信息怎样传递、迁徙怎样改变户籍和劳动、宗教与亲属网络怎样提供或限制协作。现有证据往往只照亮少数地点；保留这种不均衡，才能避免把宋辽夏金元的多政权历史改写成单一中心的必然过程。

账本条目\texttt{SS-1129-EVENT-113}所记1129年的“苗傅、刘正彦发动兵变，随后被平定”，属于南宋的历史语境。条目所示前因是前期边防、财政和统治联盟的压力构成直接背景，具体条件须按本条材料分辨，其可见后果为改变后续资源配置与政治选择；实际执行范围和社会影响仍须分项核验。这一变化若进入区域生活，必须通过道路、文书、军队、市场、家庭和地方中介才会发生；政权易手或法令发布不等于所有城镇、乡村与边地在同一时间获得同样经验。

本条所列来源为《宋史》卷24—47〈本纪〉；《建炎以来系年要录》本年条；《宋会要辑稿》相关门类。它们能够支持特定的年序、制度语言、政治行动或局部遗存，却不能无条件化为社会总量。账本特别提醒“编年主干可靠；细节、规模与动机须与对方材料、地方文书或考古材料互校”。因此，军报的兵数、条约的名分、官府的族群分类、判牍中的道德判断以及后出的叙述都须回到产生它们的场景，而不可直接代表全域动机或人口。

将事件与地方层次连接时，应同时考察资源怎样被征集、信息怎样传递、迁徙怎样改变户籍和劳动、宗教与亲属网络怎样提供或限制协作。现有证据往往只照亮少数地点；保留这种不均衡，才能避免把宋辽夏金元的多政权历史改写成单一中心的必然过程。

账本条目\texttt{SS-1129-REGNAL-YEAR}所记1129年的“南宋以宋高宗在位第3年作为诏令、赋役与外交文书的纪年框架”，属于南宋的历史语境。条目所示前因是新岁颁历、年号和纪年是中枢与地方行政沟通的共同前提，其可见后果为为同年政令、战事和地方材料提供日期锚点，不能单独推知具体政策执行。这一变化若进入区域生活，必须通过道路、文书、军队、市场、家庭和地方中介才会发生；政权易手或法令发布不等于所有城镇、乡村与边地在同一时间获得同样经验。

本条所列来源为《宋史》卷24—47〈本纪〉；《建炎以来系年要录》本年条；《宋会要辑稿》相关门类。它们能够支持特定的年序、制度语言、政治行动或局部遗存，却不能无条件化为社会总量。账本特别提醒“纪年框架可靠；该记录仅确证政权纪年连续，年内具体事项须由本年条另证”。因此，军报的兵数、条约的名分、官府的族群分类、判牍中的道德判断以及后出的叙述都须回到产生它们的场景，而不可直接代表全域动机或人口。

将事件与地方层次连接时，应同时考察资源怎样被征集、信息怎样传递、迁徙怎样改变户籍和劳动、宗教与亲属网络怎样提供或限制协作。现有证据往往只照亮少数地点；保留这种不均衡，才能避免把宋辽夏金元的多政权历史改写成单一中心的必然过程。

账本条目\texttt{SS-1130-EVENT-114}所记1130年的“金军渡江后北撤，南宋保住江南核心地区”，属于南宋与金的历史语境。条目所示前因是前期边防、财政和统治联盟的压力构成直接背景，具体条件须按本条材料分辨，其可见后果为改变后续资源配置与政治选择；实际执行范围和社会影响仍须分项核验。这一变化若进入区域生活，必须通过道路、文书、军队、市场、家庭和地方中介才会发生；政权易手或法令发布不等于所有城镇、乡村与边地在同一时间获得同样经验。

本条所列来源为《宋史》卷24—47〈本纪〉；《建炎以来系年要录》本年条；《宋会要辑稿》相关门类。它们能够支持特定的年序、制度语言、政治行动或局部遗存，却不能无条件化为社会总量。账本特别提醒“编年主干可靠；细节、规模与动机须与对方材料、地方文书或考古材料互校”。因此，军报的兵数、条约的名分、官府的族群分类、判牍中的道德判断以及后出的叙述都须回到产生它们的场景，而不可直接代表全域动机或人口。

将事件与地方层次连接时，应同时考察资源怎样被征集、信息怎样传递、迁徙怎样改变户籍和劳动、宗教与亲属网络怎样提供或限制协作。现有证据往往只照亮少数地点；保留这种不均衡，才能避免把宋辽夏金元的多政权历史改写成单一中心的必然过程。

账本条目\texttt{SS-1130-REGNAL-YEAR}所记1130年的“南宋以宋高宗在位第4年作为诏令、赋役与外交文书的纪年框架”，属于南宋的历史语境。条目所示前因是新岁颁历、年号和纪年是中枢与地方行政沟通的共同前提，其可见后果为为同年政令、战事和地方材料提供日期锚点，不能单独推知具体政策执行。这一变化若进入区域生活，必须通过道路、文书、军队、市场、家庭和地方中介才会发生；政权易手或法令发布不等于所有城镇、乡村与边地在同一时间获得同样经验。

本条所列来源为《宋史》卷24—47〈本纪〉；《建炎以来系年要录》本年条；《宋会要辑稿》相关门类。它们能够支持特定的年序、制度语言、政治行动或局部遗存，却不能无条件化为社会总量。账本特别提醒“纪年框架可靠；该记录仅确证政权纪年连续，年内具体事项须由本年条另证”。因此，军报的兵数、条约的名分、官府的族群分类、判牍中的道德判断以及后出的叙述都须回到产生它们的场景，而不可直接代表全域动机或人口。

将事件与地方层次连接时，应同时考察资源怎样被征集、信息怎样传递、迁徙怎样改变户籍和劳动、宗教与亲属网络怎样提供或限制协作。现有证据往往只照亮少数地点；保留这种不均衡，才能避免把宋辽夏金元的多政权历史改写成单一中心的必然过程。

账本条目\texttt{SS-1131-REGNAL-YEAR}所记1131年的“南宋以宋高宗在位第5年作为诏令、赋役与外交文书的纪年框架”，属于南宋的历史语境。条目所示前因是新岁颁历、年号和纪年是中枢与地方行政沟通的共同前提，其可见后果为为同年政令、战事和地方材料提供日期锚点，不能单独推知具体政策执行。这一变化若进入区域生活，必须通过道路、文书、军队、市场、家庭和地方中介才会发生；政权易手或法令发布不等于所有城镇、乡村与边地在同一时间获得同样经验。

本条所列来源为《宋史》卷24—47〈本纪〉；《建炎以来系年要录》本年条；《宋会要辑稿》相关门类。它们能够支持特定的年序、制度语言、政治行动或局部遗存，却不能无条件化为社会总量。账本特别提醒“纪年框架可靠；该记录仅确证政权纪年连续，年内具体事项须由本年条另证”。因此，军报的兵数、条约的名分、官府的族群分类、判牍中的道德判断以及后出的叙述都须回到产生它们的场景，而不可直接代表全域动机或人口。

将事件与地方层次连接时，应同时考察资源怎样被征集、信息怎样传递、迁徙怎样改变户籍和劳动、宗教与亲属网络怎样提供或限制协作。现有证据往往只照亮少数地点；保留这种不均衡，才能避免把宋辽夏金元的多政权历史改写成单一中心的必然过程。

账本条目\texttt{SS-1132-EVENT-115}所记1132年的“行在逐步以临安为稳定基地”，属于南宋的历史语境。条目所示前因是前期边防、财政和统治联盟的压力构成直接背景，具体条件须按本条材料分辨，其可见后果为改变后续资源配置与政治选择；实际执行范围和社会影响仍须分项核验。这一变化若进入区域生活，必须通过道路、文书、军队、市场、家庭和地方中介才会发生；政权易手或法令发布不等于所有城镇、乡村与边地在同一时间获得同样经验。

本条所列来源为《宋史》卷24—47〈本纪〉；《建炎以来系年要录》本年条；《宋会要辑稿》相关门类。它们能够支持特定的年序、制度语言、政治行动或局部遗存，却不能无条件化为社会总量。账本特别提醒“编年主干可靠；细节、规模与动机须与对方材料、地方文书或考古材料互校”。因此，军报的兵数、条约的名分、官府的族群分类、判牍中的道德判断以及后出的叙述都须回到产生它们的场景，而不可直接代表全域动机或人口。

将事件与地方层次连接时，应同时考察资源怎样被征集、信息怎样传递、迁徙怎样改变户籍和劳动、宗教与亲属网络怎样提供或限制协作。现有证据往往只照亮少数地点；保留这种不均衡，才能避免把宋辽夏金元的多政权历史改写成单一中心的必然过程。

账本条目\texttt{SS-1132-REGNAL-YEAR}所记1132年的“南宋以宋高宗在位第6年作为诏令、赋役与外交文书的纪年框架”，属于南宋的历史语境。条目所示前因是新岁颁历、年号和纪年是中枢与地方行政沟通的共同前提，其可见后果为为同年政令、战事和地方材料提供日期锚点，不能单独推知具体政策执行。这一变化若进入区域生活，必须通过道路、文书、军队、市场、家庭和地方中介才会发生；政权易手或法令发布不等于所有城镇、乡村与边地在同一时间获得同样经验。

本条所列来源为《宋史》卷24—47〈本纪〉；《建炎以来系年要录》本年条；《宋会要辑稿》相关门类。它们能够支持特定的年序、制度语言、政治行动或局部遗存，却不能无条件化为社会总量。账本特别提醒“纪年框架可靠；该记录仅确证政权纪年连续，年内具体事项须由本年条另证”。因此，军报的兵数、条约的名分、官府的族群分类、判牍中的道德判断以及后出的叙述都须回到产生它们的场景，而不可直接代表全域动机或人口。

将事件与地方层次连接时，应同时考察资源怎样被征集、信息怎样传递、迁徙怎样改变户籍和劳动、宗教与亲属网络怎样提供或限制协作。现有证据往往只照亮少数地点；保留这种不均衡，才能避免把宋辽夏金元的多政权历史改写成单一中心的必然过程。

账本条目\texttt{SS-1133-REGNAL-YEAR}所记1133年的“南宋以宋高宗在位第7年作为诏令、赋役与外交文书的纪年框架”，属于南宋的历史语境。条目所示前因是新岁颁历、年号和纪年是中枢与地方行政沟通的共同前提，其可见后果为为同年政令、战事和地方材料提供日期锚点，不能单独推知具体政策执行。这一变化若进入区域生活，必须通过道路、文书、军队、市场、家庭和地方中介才会发生；政权易手或法令发布不等于所有城镇、乡村与边地在同一时间获得同样经验。

本条所列来源为《宋史》卷24—47〈本纪〉；《建炎以来系年要录》本年条；《宋会要辑稿》相关门类。它们能够支持特定的年序、制度语言、政治行动或局部遗存，却不能无条件化为社会总量。账本特别提醒“纪年框架可靠；该记录仅确证政权纪年连续，年内具体事项须由本年条另证”。因此，军报的兵数、条约的名分、官府的族群分类、判牍中的道德判断以及后出的叙述都须回到产生它们的场景，而不可直接代表全域动机或人口。

将事件与地方层次连接时，应同时考察资源怎样被征集、信息怎样传递、迁徙怎样改变户籍和劳动、宗教与亲属网络怎样提供或限制协作。现有证据往往只照亮少数地点；保留这种不均衡，才能避免把宋辽夏金元的多政权历史改写成单一中心的必然过程。

\subsection{外交、道路与人口移动}

账本条目\texttt{SS-1134-EVENT-116}所记1134年的“南宋军在襄汉、淮西等方向取得局部胜利”，属于南宋与金的历史语境。条目所示前因是前期边防、财政和统治联盟的压力构成直接背景，具体条件须按本条材料分辨，其可见后果为改变后续资源配置与政治选择；实际执行范围和社会影响仍须分项核验。这一变化若进入区域生活，必须通过道路、文书、军队、市场、家庭和地方中介才会发生；政权易手或法令发布不等于所有城镇、乡村与边地在同一时间获得同样经验。

本条所列来源为《宋史》卷24—47〈本纪〉；《建炎以来系年要录》本年条；《宋会要辑稿》相关门类。它们能够支持特定的年序、制度语言、政治行动或局部遗存，却不能无条件化为社会总量。账本特别提醒“编年主干可靠；细节、规模与动机须与对方材料、地方文书或考古材料互校”。因此，军报的兵数、条约的名分、官府的族群分类、判牍中的道德判断以及后出的叙述都须回到产生它们的场景，而不可直接代表全域动机或人口。

将事件与地方层次连接时，应同时考察资源怎样被征集、信息怎样传递、迁徙怎样改变户籍和劳动、宗教与亲属网络怎样提供或限制协作。现有证据往往只照亮少数地点；保留这种不均衡，才能避免把宋辽夏金元的多政权历史改写成单一中心的必然过程。

账本条目\texttt{SS-1134-REGNAL-YEAR}所记1134年的“南宋以宋高宗在位第8年作为诏令、赋役与外交文书的纪年框架”，属于南宋的历史语境。条目所示前因是新岁颁历、年号和纪年是中枢与地方行政沟通的共同前提，其可见后果为为同年政令、战事和地方材料提供日期锚点，不能单独推知具体政策执行。这一变化若进入区域生活，必须通过道路、文书、军队、市场、家庭和地方中介才会发生；政权易手或法令发布不等于所有城镇、乡村与边地在同一时间获得同样经验。

本条所列来源为《宋史》卷24—47〈本纪〉；《建炎以来系年要录》本年条；《宋会要辑稿》相关门类。它们能够支持特定的年序、制度语言、政治行动或局部遗存，却不能无条件化为社会总量。账本特别提醒“纪年框架可靠；该记录仅确证政权纪年连续，年内具体事项须由本年条另证”。因此，军报的兵数、条约的名分、官府的族群分类、判牍中的道德判断以及后出的叙述都须回到产生它们的场景，而不可直接代表全域动机或人口。

将事件与地方层次连接时，应同时考察资源怎样被征集、信息怎样传递、迁徙怎样改变户籍和劳动、宗教与亲属网络怎样提供或限制协作。现有证据往往只照亮少数地点；保留这种不均衡，才能避免把宋辽夏金元的多政权历史改写成单一中心的必然过程。

账本条目\texttt{SS-1135-EVENT-117}所记1135年的“岳飞等将领获得更大经略与军队整编责任”，属于南宋的历史语境。条目所示前因是前期边防、财政和统治联盟的压力构成直接背景，具体条件须按本条材料分辨，其可见后果为改变后续资源配置与政治选择；实际执行范围和社会影响仍须分项核验。这一变化若进入区域生活，必须通过道路、文书、军队、市场、家庭和地方中介才会发生；政权易手或法令发布不等于所有城镇、乡村与边地在同一时间获得同样经验。

本条所列来源为《宋史》卷24—47〈本纪〉；《建炎以来系年要录》本年条；《宋会要辑稿》相关门类。它们能够支持特定的年序、制度语言、政治行动或局部遗存，却不能无条件化为社会总量。账本特别提醒“编年主干可靠；细节、规模与动机须与对方材料、地方文书或考古材料互校”。因此，军报的兵数、条约的名分、官府的族群分类、判牍中的道德判断以及后出的叙述都须回到产生它们的场景，而不可直接代表全域动机或人口。

将事件与地方层次连接时，应同时考察资源怎样被征集、信息怎样传递、迁徙怎样改变户籍和劳动、宗教与亲属网络怎样提供或限制协作。现有证据往往只照亮少数地点；保留这种不均衡，才能避免把宋辽夏金元的多政权历史改写成单一中心的必然过程。

账本条目\texttt{SS-1135-REGNAL-YEAR}所记1135年的“南宋以宋高宗在位第9年作为诏令、赋役与外交文书的纪年框架”，属于南宋的历史语境。条目所示前因是新岁颁历、年号和纪年是中枢与地方行政沟通的共同前提，其可见后果为为同年政令、战事和地方材料提供日期锚点，不能单独推知具体政策执行。这一变化若进入区域生活，必须通过道路、文书、军队、市场、家庭和地方中介才会发生；政权易手或法令发布不等于所有城镇、乡村与边地在同一时间获得同样经验。

本条所列来源为《宋史》卷24—47〈本纪〉；《建炎以来系年要录》本年条；《宋会要辑稿》相关门类。它们能够支持特定的年序、制度语言、政治行动或局部遗存，却不能无条件化为社会总量。账本特别提醒“纪年框架可靠；该记录仅确证政权纪年连续，年内具体事项须由本年条另证”。因此，军报的兵数、条约的名分、官府的族群分类、判牍中的道德判断以及后出的叙述都须回到产生它们的场景，而不可直接代表全域动机或人口。

将事件与地方层次连接时，应同时考察资源怎样被征集、信息怎样传递、迁徙怎样改变户籍和劳动、宗教与亲属网络怎样提供或限制协作。现有证据往往只照亮少数地点；保留这种不均衡，才能避免把宋辽夏金元的多政权历史改写成单一中心的必然过程。

账本条目\texttt{SS-1136-REGNAL-YEAR}所记1136年的“南宋以宋高宗在位第10年作为诏令、赋役与外交文书的纪年框架”，属于南宋的历史语境。条目所示前因是新岁颁历、年号和纪年是中枢与地方行政沟通的共同前提，其可见后果为为同年政令、战事和地方材料提供日期锚点，不能单独推知具体政策执行。这一变化若进入区域生活，必须通过道路、文书、军队、市场、家庭和地方中介才会发生；政权易手或法令发布不等于所有城镇、乡村与边地在同一时间获得同样经验。

本条所列来源为《宋史》卷24—47〈本纪〉；《建炎以来系年要录》本年条；《宋会要辑稿》相关门类。它们能够支持特定的年序、制度语言、政治行动或局部遗存，却不能无条件化为社会总量。账本特别提醒“纪年框架可靠；该记录仅确证政权纪年连续，年内具体事项须由本年条另证”。因此，军报的兵数、条约的名分、官府的族群分类、判牍中的道德判断以及后出的叙述都须回到产生它们的场景，而不可直接代表全域动机或人口。

将事件与地方层次连接时，应同时考察资源怎样被征集、信息怎样传递、迁徙怎样改变户籍和劳动、宗教与亲属网络怎样提供或限制协作。现有证据往往只照亮少数地点；保留这种不均衡，才能避免把宋辽夏金元的多政权历史改写成单一中心的必然过程。

账本条目\texttt{SS-1137-REGNAL-YEAR}所记1137年的“南宋以宋高宗在位第11年作为诏令、赋役与外交文书的纪年框架”，属于南宋的历史语境。条目所示前因是新岁颁历、年号和纪年是中枢与地方行政沟通的共同前提，其可见后果为为同年政令、战事和地方材料提供日期锚点，不能单独推知具体政策执行。这一变化若进入区域生活，必须通过道路、文书、军队、市场、家庭和地方中介才会发生；政权易手或法令发布不等于所有城镇、乡村与边地在同一时间获得同样经验。

本条所列来源为《宋史》卷24—47〈本纪〉；《建炎以来系年要录》本年条；《宋会要辑稿》相关门类。它们能够支持特定的年序、制度语言、政治行动或局部遗存，却不能无条件化为社会总量。账本特别提醒“纪年框架可靠；该记录仅确证政权纪年连续，年内具体事项须由本年条另证”。因此，军报的兵数、条约的名分、官府的族群分类、判牍中的道德判断以及后出的叙述都须回到产生它们的场景，而不可直接代表全域动机或人口。

将事件与地方层次连接时，应同时考察资源怎样被征集、信息怎样传递、迁徙怎样改变户籍和劳动、宗教与亲属网络怎样提供或限制协作。现有证据往往只照亮少数地点；保留这种不均衡，才能避免把宋辽夏金元的多政权历史改写成单一中心的必然过程。

账本条目\texttt{SS-1138-EVENT-118}所记1138年的“临安被确立为行在和事实都城”，属于南宋的历史语境。条目所示前因是前期边防、财政和统治联盟的压力构成直接背景，具体条件须按本条材料分辨，其可见后果为改变后续资源配置与政治选择；实际执行范围和社会影响仍须分项核验。这一变化若进入区域生活，必须通过道路、文书、军队、市场、家庭和地方中介才会发生；政权易手或法令发布不等于所有城镇、乡村与边地在同一时间获得同样经验。

本条所列来源为《宋史》卷24—47〈本纪〉；《建炎以来系年要录》本年条；《宋会要辑稿》相关门类。它们能够支持特定的年序、制度语言、政治行动或局部遗存，却不能无条件化为社会总量。账本特别提醒“编年主干可靠；细节、规模与动机须与对方材料、地方文书或考古材料互校”。因此，军报的兵数、条约的名分、官府的族群分类、判牍中的道德判断以及后出的叙述都须回到产生它们的场景，而不可直接代表全域动机或人口。

将事件与地方层次连接时，应同时考察资源怎样被征集、信息怎样传递、迁徙怎样改变户籍和劳动、宗教与亲属网络怎样提供或限制协作。现有证据往往只照亮少数地点；保留这种不均衡，才能避免把宋辽夏金元的多政权历史改写成单一中心的必然过程。

账本条目\texttt{SS-1138-REGNAL-YEAR}所记1138年的“南宋以宋高宗在位第12年作为诏令、赋役与外交文书的纪年框架”，属于南宋的历史语境。条目所示前因是新岁颁历、年号和纪年是中枢与地方行政沟通的共同前提，其可见后果为为同年政令、战事和地方材料提供日期锚点，不能单独推知具体政策执行。这一变化若进入区域生活，必须通过道路、文书、军队、市场、家庭和地方中介才会发生；政权易手或法令发布不等于所有城镇、乡村与边地在同一时间获得同样经验。

本条所列来源为《宋史》卷24—47〈本纪〉；《建炎以来系年要录》本年条；《宋会要辑稿》相关门类。它们能够支持特定的年序、制度语言、政治行动或局部遗存，却不能无条件化为社会总量。账本特别提醒“纪年框架可靠；该记录仅确证政权纪年连续，年内具体事项须由本年条另证”。因此，军报的兵数、条约的名分、官府的族群分类、判牍中的道德判断以及后出的叙述都须回到产生它们的场景，而不可直接代表全域动机或人口。

将事件与地方层次连接时，应同时考察资源怎样被征集、信息怎样传递、迁徙怎样改变户籍和劳动、宗教与亲属网络怎样提供或限制协作。现有证据往往只照亮少数地点；保留这种不均衡，才能避免把宋辽夏金元的多政权历史改写成单一中心的必然过程。

账本条目\texttt{SS-1139-REGNAL-YEAR}所记1139年的“南宋以宋高宗在位第13年作为诏令、赋役与外交文书的纪年框架”，属于南宋的历史语境。条目所示前因是新岁颁历、年号和纪年是中枢与地方行政沟通的共同前提，其可见后果为为同年政令、战事和地方材料提供日期锚点，不能单独推知具体政策执行。这一变化若进入区域生活，必须通过道路、文书、军队、市场、家庭和地方中介才会发生；政权易手或法令发布不等于所有城镇、乡村与边地在同一时间获得同样经验。

本条所列来源为《宋史》卷24—47〈本纪〉；《建炎以来系年要录》本年条；《宋会要辑稿》相关门类。它们能够支持特定的年序、制度语言、政治行动或局部遗存，却不能无条件化为社会总量。账本特别提醒“纪年框架可靠；该记录仅确证政权纪年连续，年内具体事项须由本年条另证”。因此，军报的兵数、条约的名分、官府的族群分类、判牍中的道德判断以及后出的叙述都须回到产生它们的场景，而不可直接代表全域动机或人口。

将事件与地方层次连接时，应同时考察资源怎样被征集、信息怎样传递、迁徙怎样改变户籍和劳动、宗教与亲属网络怎样提供或限制协作。现有证据往往只照亮少数地点；保留这种不均衡，才能避免把宋辽夏金元的多政权历史改写成单一中心的必然过程。

账本条目\texttt{SS-1140-EVENT-119}所记1140年的“南宋军在河南、淮西等方向反攻”，属于南宋与金的历史语境。条目所示前因是前期边防、财政和统治联盟的压力构成直接背景，具体条件须按本条材料分辨，其可见后果为改变后续资源配置与政治选择；实际执行范围和社会影响仍须分项核验。这一变化若进入区域生活，必须通过道路、文书、军队、市场、家庭和地方中介才会发生；政权易手或法令发布不等于所有城镇、乡村与边地在同一时间获得同样经验。

本条所列来源为《宋史》卷24—47〈本纪〉；《建炎以来系年要录》本年条；《宋会要辑稿》相关门类。它们能够支持特定的年序、制度语言、政治行动或局部遗存，却不能无条件化为社会总量。账本特别提醒“编年主干可靠；细节、规模与动机须与对方材料、地方文书或考古材料互校”。因此，军报的兵数、条约的名分、官府的族群分类、判牍中的道德判断以及后出的叙述都须回到产生它们的场景，而不可直接代表全域动机或人口。

将事件与地方层次连接时，应同时考察资源怎样被征集、信息怎样传递、迁徙怎样改变户籍和劳动、宗教与亲属网络怎样提供或限制协作。现有证据往往只照亮少数地点；保留这种不均衡，才能避免把宋辽夏金元的多政权历史改写成单一中心的必然过程。

账本条目\texttt{SS-1140-REGNAL-YEAR}所记1140年的“南宋以宋高宗在位第14年作为诏令、赋役与外交文书的纪年框架”，属于南宋的历史语境。条目所示前因是新岁颁历、年号和纪年是中枢与地方行政沟通的共同前提，其可见后果为为同年政令、战事和地方材料提供日期锚点，不能单独推知具体政策执行。这一变化若进入区域生活，必须通过道路、文书、军队、市场、家庭和地方中介才会发生；政权易手或法令发布不等于所有城镇、乡村与边地在同一时间获得同样经验。

本条所列来源为《宋史》卷24—47〈本纪〉；《建炎以来系年要录》本年条；《宋会要辑稿》相关门类。它们能够支持特定的年序、制度语言、政治行动或局部遗存，却不能无条件化为社会总量。账本特别提醒“纪年框架可靠；该记录仅确证政权纪年连续，年内具体事项须由本年条另证”。因此，军报的兵数、条约的名分、官府的族群分类、判牍中的道德判断以及后出的叙述都须回到产生它们的场景，而不可直接代表全域动机或人口。

将事件与地方层次连接时，应同时考察资源怎样被征集、信息怎样传递、迁徙怎样改变户籍和劳动、宗教与亲属网络怎样提供或限制协作。现有证据往往只照亮少数地点；保留这种不均衡，才能避免把宋辽夏金元的多政权历史改写成单一中心的必然过程。

账本条目\texttt{SS-1141-EVENT-120}所记1141年的“南宋接受绍兴和议框架”，属于南宋与金的历史语境。条目所示前因是前期边防、财政和统治联盟的压力构成直接背景，具体条件须按本条材料分辨，其可见后果为改变后续资源配置与政治选择；实际执行范围和社会影响仍须分项核验。这一变化若进入区域生活，必须通过道路、文书、军队、市场、家庭和地方中介才会发生；政权易手或法令发布不等于所有城镇、乡村与边地在同一时间获得同样经验。

本条所列来源为《宋史》卷24—47〈本纪〉；《建炎以来系年要录》本年条；《宋会要辑稿》相关门类。它们能够支持特定的年序、制度语言、政治行动或局部遗存，却不能无条件化为社会总量。账本特别提醒“编年主干可靠；细节、规模与动机须与对方材料、地方文书或考古材料互校”。因此，军报的兵数、条约的名分、官府的族群分类、判牍中的道德判断以及后出的叙述都须回到产生它们的场景，而不可直接代表全域动机或人口。

将事件与地方层次连接时，应同时考察资源怎样被征集、信息怎样传递、迁徙怎样改变户籍和劳动、宗教与亲属网络怎样提供或限制协作。现有证据往往只照亮少数地点；保留这种不均衡，才能避免把宋辽夏金元的多政权历史改写成单一中心的必然过程。

\subsection{环境风险与地方经济}

账本条目\texttt{SS-1141-REGNAL-YEAR}所记1141年的“南宋以宋高宗在位第15年作为诏令、赋役与外交文书的纪年框架”，属于南宋的历史语境。条目所示前因是新岁颁历、年号和纪年是中枢与地方行政沟通的共同前提，其可见后果为为同年政令、战事和地方材料提供日期锚点，不能单独推知具体政策执行。这一变化若进入区域生活，必须通过道路、文书、军队、市场、家庭和地方中介才会发生；政权易手或法令发布不等于所有城镇、乡村与边地在同一时间获得同样经验。

本条所列来源为《宋史》卷24—47〈本纪〉；《建炎以来系年要录》本年条；《宋会要辑稿》相关门类。它们能够支持特定的年序、制度语言、政治行动或局部遗存，却不能无条件化为社会总量。账本特别提醒“纪年框架可靠；该记录仅确证政权纪年连续，年内具体事项须由本年条另证”。因此，军报的兵数、条约的名分、官府的族群分类、判牍中的道德判断以及后出的叙述都须回到产生它们的场景，而不可直接代表全域动机或人口。

将事件与地方层次连接时，应同时考察资源怎样被征集、信息怎样传递、迁徙怎样改变户籍和劳动、宗教与亲属网络怎样提供或限制协作。现有证据往往只照亮少数地点；保留这种不均衡，才能避免把宋辽夏金元的多政权历史改写成单一中心的必然过程。

账本条目\texttt{SS-1142-EVENT-121}所记1142年的“岳飞被处死，相关将领和军队受到整顿”，属于南宋的历史语境。条目所示前因是前期边防、财政和统治联盟的压力构成直接背景，具体条件须按本条材料分辨，其可见后果为改变后续资源配置与政治选择；实际执行范围和社会影响仍须分项核验。这一变化若进入区域生活，必须通过道路、文书、军队、市场、家庭和地方中介才会发生；政权易手或法令发布不等于所有城镇、乡村与边地在同一时间获得同样经验。

本条所列来源为《宋史》卷24—47〈本纪〉；《建炎以来系年要录》本年条；《宋会要辑稿》相关门类。它们能够支持特定的年序、制度语言、政治行动或局部遗存，却不能无条件化为社会总量。账本特别提醒“编年主干可靠；细节、规模与动机须与对方材料、地方文书或考古材料互校”。因此，军报的兵数、条约的名分、官府的族群分类、判牍中的道德判断以及后出的叙述都须回到产生它们的场景，而不可直接代表全域动机或人口。

将事件与地方层次连接时，应同时考察资源怎样被征集、信息怎样传递、迁徙怎样改变户籍和劳动、宗教与亲属网络怎样提供或限制协作。现有证据往往只照亮少数地点；保留这种不均衡，才能避免把宋辽夏金元的多政权历史改写成单一中心的必然过程。

账本条目\texttt{SS-1142-REGNAL-YEAR}所记1142年的“南宋以宋高宗在位第16年作为诏令、赋役与外交文书的纪年框架”，属于南宋的历史语境。条目所示前因是新岁颁历、年号和纪年是中枢与地方行政沟通的共同前提，其可见后果为为同年政令、战事和地方材料提供日期锚点，不能单独推知具体政策执行。这一变化若进入区域生活，必须通过道路、文书、军队、市场、家庭和地方中介才会发生；政权易手或法令发布不等于所有城镇、乡村与边地在同一时间获得同样经验。

本条所列来源为《宋史》卷24—47〈本纪〉；《建炎以来系年要录》本年条；《宋会要辑稿》相关门类。它们能够支持特定的年序、制度语言、政治行动或局部遗存，却不能无条件化为社会总量。账本特别提醒“纪年框架可靠；该记录仅确证政权纪年连续，年内具体事项须由本年条另证”。因此，军报的兵数、条约的名分、官府的族群分类、判牍中的道德判断以及后出的叙述都须回到产生它们的场景，而不可直接代表全域动机或人口。

将事件与地方层次连接时，应同时考察资源怎样被征集、信息怎样传递、迁徙怎样改变户籍和劳动、宗教与亲属网络怎样提供或限制协作。现有证据往往只照亮少数地点；保留这种不均衡，才能避免把宋辽夏金元的多政权历史改写成单一中心的必然过程。

账本条目\texttt{SS-1143-REGNAL-YEAR}所记1143年的“南宋以宋高宗在位第17年作为诏令、赋役与外交文书的纪年框架”，属于南宋的历史语境。条目所示前因是新岁颁历、年号和纪年是中枢与地方行政沟通的共同前提，其可见后果为为同年政令、战事和地方材料提供日期锚点，不能单独推知具体政策执行。这一变化若进入区域生活，必须通过道路、文书、军队、市场、家庭和地方中介才会发生；政权易手或法令发布不等于所有城镇、乡村与边地在同一时间获得同样经验。

本条所列来源为《宋史》卷24—47〈本纪〉；《建炎以来系年要录》本年条；《宋会要辑稿》相关门类。它们能够支持特定的年序、制度语言、政治行动或局部遗存，却不能无条件化为社会总量。账本特别提醒“纪年框架可靠；该记录仅确证政权纪年连续，年内具体事项须由本年条另证”。因此，军报的兵数、条约的名分、官府的族群分类、判牍中的道德判断以及后出的叙述都须回到产生它们的场景，而不可直接代表全域动机或人口。

将事件与地方层次连接时，应同时考察资源怎样被征集、信息怎样传递、迁徙怎样改变户籍和劳动、宗教与亲属网络怎样提供或限制协作。现有证据往往只照亮少数地点；保留这种不均衡，才能避免把宋辽夏金元的多政权历史改写成单一中心的必然过程。

账本条目\texttt{SS-1144-REGNAL-YEAR}所记1144年的“南宋以宋高宗在位第18年作为诏令、赋役与外交文书的纪年框架”，属于南宋的历史语境。条目所示前因是新岁颁历、年号和纪年是中枢与地方行政沟通的共同前提，其可见后果为为同年政令、战事和地方材料提供日期锚点，不能单独推知具体政策执行。这一变化若进入区域生活，必须通过道路、文书、军队、市场、家庭和地方中介才会发生；政权易手或法令发布不等于所有城镇、乡村与边地在同一时间获得同样经验。

本条所列来源为《宋史》卷24—47〈本纪〉；《建炎以来系年要录》本年条；《宋会要辑稿》相关门类。它们能够支持特定的年序、制度语言、政治行动或局部遗存，却不能无条件化为社会总量。账本特别提醒“纪年框架可靠；该记录仅确证政权纪年连续，年内具体事项须由本年条另证”。因此，军报的兵数、条约的名分、官府的族群分类、判牍中的道德判断以及后出的叙述都须回到产生它们的场景，而不可直接代表全域动机或人口。

将事件与地方层次连接时，应同时考察资源怎样被征集、信息怎样传递、迁徙怎样改变户籍和劳动、宗教与亲属网络怎样提供或限制协作。现有证据往往只照亮少数地点；保留这种不均衡，才能避免把宋辽夏金元的多政权历史改写成单一中心的必然过程。

账本条目\texttt{SS-1145-REGNAL-YEAR}所记1145年的“南宋以宋高宗在位第19年作为诏令、赋役与外交文书的纪年框架”，属于南宋的历史语境。条目所示前因是新岁颁历、年号和纪年是中枢与地方行政沟通的共同前提，其可见后果为为同年政令、战事和地方材料提供日期锚点，不能单独推知具体政策执行。这一变化若进入区域生活，必须通过道路、文书、军队、市场、家庭和地方中介才会发生；政权易手或法令发布不等于所有城镇、乡村与边地在同一时间获得同样经验。

本条所列来源为《宋史》卷24—47〈本纪〉；《建炎以来系年要录》本年条；《宋会要辑稿》相关门类。它们能够支持特定的年序、制度语言、政治行动或局部遗存，却不能无条件化为社会总量。账本特别提醒“纪年框架可靠；该记录仅确证政权纪年连续，年内具体事项须由本年条另证”。因此，军报的兵数、条约的名分、官府的族群分类、判牍中的道德判断以及后出的叙述都须回到产生它们的场景，而不可直接代表全域动机或人口。

将事件与地方层次连接时，应同时考察资源怎样被征集、信息怎样传递、迁徙怎样改变户籍和劳动、宗教与亲属网络怎样提供或限制协作。现有证据往往只照亮少数地点；保留这种不均衡，才能避免把宋辽夏金元的多政权历史改写成单一中心的必然过程。

账本条目\texttt{SS-1146-REGNAL-YEAR}所记1146年的“南宋以宋高宗在位第20年作为诏令、赋役与外交文书的纪年框架”，属于南宋的历史语境。条目所示前因是新岁颁历、年号和纪年是中枢与地方行政沟通的共同前提，其可见后果为为同年政令、战事和地方材料提供日期锚点，不能单独推知具体政策执行。这一变化若进入区域生活，必须通过道路、文书、军队、市场、家庭和地方中介才会发生；政权易手或法令发布不等于所有城镇、乡村与边地在同一时间获得同样经验。

本条所列来源为《宋史》卷24—47〈本纪〉；《建炎以来系年要录》本年条；《宋会要辑稿》相关门类。它们能够支持特定的年序、制度语言、政治行动或局部遗存，却不能无条件化为社会总量。账本特别提醒“纪年框架可靠；该记录仅确证政权纪年连续，年内具体事项须由本年条另证”。因此，军报的兵数、条约的名分、官府的族群分类、判牍中的道德判断以及后出的叙述都须回到产生它们的场景，而不可直接代表全域动机或人口。

将事件与地方层次连接时，应同时考察资源怎样被征集、信息怎样传递、迁徙怎样改变户籍和劳动、宗教与亲属网络怎样提供或限制协作。现有证据往往只照亮少数地点；保留这种不均衡，才能避免把宋辽夏金元的多政权历史改写成单一中心的必然过程。

账本条目\texttt{SS-1147-REGNAL-YEAR}所记1147年的“南宋以宋高宗在位第21年作为诏令、赋役与外交文书的纪年框架”，属于南宋的历史语境。条目所示前因是新岁颁历、年号和纪年是中枢与地方行政沟通的共同前提，其可见后果为为同年政令、战事和地方材料提供日期锚点，不能单独推知具体政策执行。这一变化若进入区域生活，必须通过道路、文书、军队、市场、家庭和地方中介才会发生；政权易手或法令发布不等于所有城镇、乡村与边地在同一时间获得同样经验。

本条所列来源为《宋史》卷24—47〈本纪〉；《建炎以来系年要录》本年条；《宋会要辑稿》相关门类。它们能够支持特定的年序、制度语言、政治行动或局部遗存，却不能无条件化为社会总量。账本特别提醒“纪年框架可靠；该记录仅确证政权纪年连续，年内具体事项须由本年条另证”。因此，军报的兵数、条约的名分、官府的族群分类、判牍中的道德判断以及后出的叙述都须回到产生它们的场景，而不可直接代表全域动机或人口。

将事件与地方层次连接时，应同时考察资源怎样被征集、信息怎样传递、迁徙怎样改变户籍和劳动、宗教与亲属网络怎样提供或限制协作。现有证据往往只照亮少数地点；保留这种不均衡，才能避免把宋辽夏金元的多政权历史改写成单一中心的必然过程。

账本条目\texttt{SS-1148-REGNAL-YEAR}所记1148年的“南宋以宋高宗在位第22年作为诏令、赋役与外交文书的纪年框架”，属于南宋的历史语境。条目所示前因是新岁颁历、年号和纪年是中枢与地方行政沟通的共同前提，其可见后果为为同年政令、战事和地方材料提供日期锚点，不能单独推知具体政策执行。这一变化若进入区域生活，必须通过道路、文书、军队、市场、家庭和地方中介才会发生；政权易手或法令发布不等于所有城镇、乡村与边地在同一时间获得同样经验。

本条所列来源为《宋史》卷24—47〈本纪〉；《建炎以来系年要录》本年条；《宋会要辑稿》相关门类。它们能够支持特定的年序、制度语言、政治行动或局部遗存，却不能无条件化为社会总量。账本特别提醒“纪年框架可靠；该记录仅确证政权纪年连续，年内具体事项须由本年条另证”。因此，军报的兵数、条约的名分、官府的族群分类、判牍中的道德判断以及后出的叙述都须回到产生它们的场景，而不可直接代表全域动机或人口。

将事件与地方层次连接时，应同时考察资源怎样被征集、信息怎样传递、迁徙怎样改变户籍和劳动、宗教与亲属网络怎样提供或限制协作。现有证据往往只照亮少数地点；保留这种不均衡，才能避免把宋辽夏金元的多政权历史改写成单一中心的必然过程。

账本条目\texttt{SS-1149-REGNAL-YEAR}所记1149年的“南宋以宋高宗在位第23年作为诏令、赋役与外交文书的纪年框架”，属于南宋的历史语境。条目所示前因是新岁颁历、年号和纪年是中枢与地方行政沟通的共同前提，其可见后果为为同年政令、战事和地方材料提供日期锚点，不能单独推知具体政策执行。这一变化若进入区域生活，必须通过道路、文书、军队、市场、家庭和地方中介才会发生；政权易手或法令发布不等于所有城镇、乡村与边地在同一时间获得同样经验。

本条所列来源为《宋史》卷24—47〈本纪〉；《建炎以来系年要录》本年条；《宋会要辑稿》相关门类。它们能够支持特定的年序、制度语言、政治行动或局部遗存，却不能无条件化为社会总量。账本特别提醒“纪年框架可靠；该记录仅确证政权纪年连续，年内具体事项须由本年条另证”。因此，军报的兵数、条约的名分、官府的族群分类、判牍中的道德判断以及后出的叙述都须回到产生它们的场景，而不可直接代表全域动机或人口。

将事件与地方层次连接时，应同时考察资源怎样被征集、信息怎样传递、迁徙怎样改变户籍和劳动、宗教与亲属网络怎样提供或限制协作。现有证据往往只照亮少数地点；保留这种不均衡，才能避免把宋辽夏金元的多政权历史改写成单一中心的必然过程。

账本条目\texttt{SS-1150-EVENT-122}所记1150年的“朝廷调整会子与纸币管理”，属于南宋的历史语境。条目所示前因是前期边防、财政和统治联盟的压力构成直接背景，具体条件须按本条材料分辨，其可见后果为改变后续资源配置与政治选择；实际执行范围和社会影响仍须分项核验。这一变化若进入区域生活，必须通过道路、文书、军队、市场、家庭和地方中介才会发生；政权易手或法令发布不等于所有城镇、乡村与边地在同一时间获得同样经验。

本条所列来源为《宋史》卷24—47〈本纪〉；《建炎以来系年要录》本年条；《宋会要辑稿》相关门类。它们能够支持特定的年序、制度语言、政治行动或局部遗存，却不能无条件化为社会总量。账本特别提醒“编年主干可靠；细节、规模与动机须与对方材料、地方文书或考古材料互校”。因此，军报的兵数、条约的名分、官府的族群分类、判牍中的道德判断以及后出的叙述都须回到产生它们的场景，而不可直接代表全域动机或人口。

将事件与地方层次连接时，应同时考察资源怎样被征集、信息怎样传递、迁徙怎样改变户籍和劳动、宗教与亲属网络怎样提供或限制协作。现有证据往往只照亮少数地点；保留这种不均衡，才能避免把宋辽夏金元的多政权历史改写成单一中心的必然过程。

账本条目\texttt{SS-1150-REGNAL-YEAR}所记1150年的“南宋以宋高宗在位第24年作为诏令、赋役与外交文书的纪年框架”，属于南宋的历史语境。条目所示前因是新岁颁历、年号和纪年是中枢与地方行政沟通的共同前提，其可见后果为为同年政令、战事和地方材料提供日期锚点，不能单独推知具体政策执行。这一变化若进入区域生活，必须通过道路、文书、军队、市场、家庭和地方中介才会发生；政权易手或法令发布不等于所有城镇、乡村与边地在同一时间获得同样经验。

本条所列来源为《宋史》卷24—47〈本纪〉；《建炎以来系年要录》本年条；《宋会要辑稿》相关门类。它们能够支持特定的年序、制度语言、政治行动或局部遗存，却不能无条件化为社会总量。账本特别提醒“纪年框架可靠；该记录仅确证政权纪年连续，年内具体事项须由本年条另证”。因此，军报的兵数、条约的名分、官府的族群分类、判牍中的道德判断以及后出的叙述都须回到产生它们的场景，而不可直接代表全域动机或人口。

将事件与地方层次连接时，应同时考察资源怎样被征集、信息怎样传递、迁徙怎样改变户籍和劳动、宗教与亲属网络怎样提供或限制协作。现有证据往往只照亮少数地点；保留这种不均衡，才能避免把宋辽夏金元的多政权历史改写成单一中心的必然过程。

\subsection{宗教、法律与材料边界}

账本条目\texttt{SS-1151-REGNAL-YEAR}所记1151年的“南宋以宋高宗在位第25年作为诏令、赋役与外交文书的纪年框架”，属于南宋的历史语境。条目所示前因是新岁颁历、年号和纪年是中枢与地方行政沟通的共同前提，其可见后果为为同年政令、战事和地方材料提供日期锚点，不能单独推知具体政策执行。这一变化若进入区域生活，必须通过道路、文书、军队、市场、家庭和地方中介才会发生；政权易手或法令发布不等于所有城镇、乡村与边地在同一时间获得同样经验。

本条所列来源为《宋史》卷24—47〈本纪〉；《建炎以来系年要录》本年条；《宋会要辑稿》相关门类。它们能够支持特定的年序、制度语言、政治行动或局部遗存，却不能无条件化为社会总量。账本特别提醒“纪年框架可靠；该记录仅确证政权纪年连续，年内具体事项须由本年条另证”。因此，军报的兵数、条约的名分、官府的族群分类、判牍中的道德判断以及后出的叙述都须回到产生它们的场景，而不可直接代表全域动机或人口。

将事件与地方层次连接时，应同时考察资源怎样被征集、信息怎样传递、迁徙怎样改变户籍和劳动、宗教与亲属网络怎样提供或限制协作。现有证据往往只照亮少数地点；保留这种不均衡，才能避免把宋辽夏金元的多政权历史改写成单一中心的必然过程。

账本条目\texttt{SS-1152-REGNAL-YEAR}所记1152年的“南宋以宋高宗在位第26年作为诏令、赋役与外交文书的纪年框架”，属于南宋的历史语境。条目所示前因是新岁颁历、年号和纪年是中枢与地方行政沟通的共同前提，其可见后果为为同年政令、战事和地方材料提供日期锚点，不能单独推知具体政策执行。这一变化若进入区域生活，必须通过道路、文书、军队、市场、家庭和地方中介才会发生；政权易手或法令发布不等于所有城镇、乡村与边地在同一时间获得同样经验。

本条所列来源为《宋史》卷24—47〈本纪〉；《建炎以来系年要录》本年条；《宋会要辑稿》相关门类。它们能够支持特定的年序、制度语言、政治行动或局部遗存，却不能无条件化为社会总量。账本特别提醒“纪年框架可靠；该记录仅确证政权纪年连续，年内具体事项须由本年条另证”。因此，军报的兵数、条约的名分、官府的族群分类、判牍中的道德判断以及后出的叙述都须回到产生它们的场景，而不可直接代表全域动机或人口。

将事件与地方层次连接时，应同时考察资源怎样被征集、信息怎样传递、迁徙怎样改变户籍和劳动、宗教与亲属网络怎样提供或限制协作。现有证据往往只照亮少数地点；保留这种不均衡，才能避免把宋辽夏金元的多政权历史改写成单一中心的必然过程。

账本条目\texttt{SS-1153-REGNAL-YEAR}所记1153年的“南宋以宋高宗在位第27年作为诏令、赋役与外交文书的纪年框架”，属于南宋的历史语境。条目所示前因是新岁颁历、年号和纪年是中枢与地方行政沟通的共同前提，其可见后果为为同年政令、战事和地方材料提供日期锚点，不能单独推知具体政策执行。这一变化若进入区域生活，必须通过道路、文书、军队、市场、家庭和地方中介才会发生；政权易手或法令发布不等于所有城镇、乡村与边地在同一时间获得同样经验。

本条所列来源为《宋史》卷24—47〈本纪〉；《建炎以来系年要录》本年条；《宋会要辑稿》相关门类。它们能够支持特定的年序、制度语言、政治行动或局部遗存，却不能无条件化为社会总量。账本特别提醒“纪年框架可靠；该记录仅确证政权纪年连续，年内具体事项须由本年条另证”。因此，军报的兵数、条约的名分、官府的族群分类、判牍中的道德判断以及后出的叙述都须回到产生它们的场景，而不可直接代表全域动机或人口。

将事件与地方层次连接时，应同时考察资源怎样被征集、信息怎样传递、迁徙怎样改变户籍和劳动、宗教与亲属网络怎样提供或限制协作。现有证据往往只照亮少数地点；保留这种不均衡，才能避免把宋辽夏金元的多政权历史改写成单一中心的必然过程。

账本条目\texttt{SS-1154-REGNAL-YEAR}所记1154年的“南宋以宋高宗在位第28年作为诏令、赋役与外交文书的纪年框架”，属于南宋的历史语境。条目所示前因是新岁颁历、年号和纪年是中枢与地方行政沟通的共同前提，其可见后果为为同年政令、战事和地方材料提供日期锚点，不能单独推知具体政策执行。这一变化若进入区域生活，必须通过道路、文书、军队、市场、家庭和地方中介才会发生；政权易手或法令发布不等于所有城镇、乡村与边地在同一时间获得同样经验。

本条所列来源为《宋史》卷24—47〈本纪〉；《建炎以来系年要录》本年条；《宋会要辑稿》相关门类。它们能够支持特定的年序、制度语言、政治行动或局部遗存，却不能无条件化为社会总量。账本特别提醒“纪年框架可靠；该记录仅确证政权纪年连续，年内具体事项须由本年条另证”。因此，军报的兵数、条约的名分、官府的族群分类、判牍中的道德判断以及后出的叙述都须回到产生它们的场景，而不可直接代表全域动机或人口。

将事件与地方层次连接时，应同时考察资源怎样被征集、信息怎样传递、迁徙怎样改变户籍和劳动、宗教与亲属网络怎样提供或限制协作。现有证据往往只照亮少数地点；保留这种不均衡，才能避免把宋辽夏金元的多政权历史改写成单一中心的必然过程。

账本条目\texttt{SS-1155-REGNAL-YEAR}所记1155年的“南宋以宋高宗在位第29年作为诏令、赋役与外交文书的纪年框架”，属于南宋的历史语境。条目所示前因是新岁颁历、年号和纪年是中枢与地方行政沟通的共同前提，其可见后果为为同年政令、战事和地方材料提供日期锚点，不能单独推知具体政策执行。这一变化若进入区域生活，必须通过道路、文书、军队、市场、家庭和地方中介才会发生；政权易手或法令发布不等于所有城镇、乡村与边地在同一时间获得同样经验。

本条所列来源为《宋史》卷24—47〈本纪〉；《建炎以来系年要录》本年条；《宋会要辑稿》相关门类。它们能够支持特定的年序、制度语言、政治行动或局部遗存，却不能无条件化为社会总量。账本特别提醒“纪年框架可靠；该记录仅确证政权纪年连续，年内具体事项须由本年条另证”。因此，军报的兵数、条约的名分、官府的族群分类、判牍中的道德判断以及后出的叙述都须回到产生它们的场景，而不可直接代表全域动机或人口。

将事件与地方层次连接时，应同时考察资源怎样被征集、信息怎样传递、迁徙怎样改变户籍和劳动、宗教与亲属网络怎样提供或限制协作。现有证据往往只照亮少数地点；保留这种不均衡，才能避免把宋辽夏金元的多政权历史改写成单一中心的必然过程。

账本条目\texttt{SS-1156-REGNAL-YEAR}所记1156年的“南宋以宋高宗在位第30年作为诏令、赋役与外交文书的纪年框架”，属于南宋的历史语境。条目所示前因是新岁颁历、年号和纪年是中枢与地方行政沟通的共同前提，其可见后果为为同年政令、战事和地方材料提供日期锚点，不能单独推知具体政策执行。这一变化若进入区域生活，必须通过道路、文书、军队、市场、家庭和地方中介才会发生；政权易手或法令发布不等于所有城镇、乡村与边地在同一时间获得同样经验。

本条所列来源为《宋史》卷24—47〈本纪〉；《建炎以来系年要录》本年条；《宋会要辑稿》相关门类。它们能够支持特定的年序、制度语言、政治行动或局部遗存，却不能无条件化为社会总量。账本特别提醒“纪年框架可靠；该记录仅确证政权纪年连续，年内具体事项须由本年条另证”。因此，军报的兵数、条约的名分、官府的族群分类、判牍中的道德判断以及后出的叙述都须回到产生它们的场景，而不可直接代表全域动机或人口。

将事件与地方层次连接时，应同时考察资源怎样被征集、信息怎样传递、迁徙怎样改变户籍和劳动、宗教与亲属网络怎样提供或限制协作。现有证据往往只照亮少数地点；保留这种不均衡，才能避免把宋辽夏金元的多政权历史改写成单一中心的必然过程。

账本条目\texttt{SS-1157-REGNAL-YEAR}所记1157年的“南宋以宋高宗在位第31年作为诏令、赋役与外交文书的纪年框架”，属于南宋的历史语境。条目所示前因是新岁颁历、年号和纪年是中枢与地方行政沟通的共同前提，其可见后果为为同年政令、战事和地方材料提供日期锚点，不能单独推知具体政策执行。这一变化若进入区域生活，必须通过道路、文书、军队、市场、家庭和地方中介才会发生；政权易手或法令发布不等于所有城镇、乡村与边地在同一时间获得同样经验。

本条所列来源为《宋史》卷24—47〈本纪〉；《建炎以来系年要录》本年条；《宋会要辑稿》相关门类。它们能够支持特定的年序、制度语言、政治行动或局部遗存，却不能无条件化为社会总量。账本特别提醒“纪年框架可靠；该记录仅确证政权纪年连续，年内具体事项须由本年条另证”。因此，军报的兵数、条约的名分、官府的族群分类、判牍中的道德判断以及后出的叙述都须回到产生它们的场景，而不可直接代表全域动机或人口。

将事件与地方层次连接时，应同时考察资源怎样被征集、信息怎样传递、迁徙怎样改变户籍和劳动、宗教与亲属网络怎样提供或限制协作。现有证据往往只照亮少数地点；保留这种不均衡，才能避免把宋辽夏金元的多政权历史改写成单一中心的必然过程。

账本条目\texttt{SS-1158-REGNAL-YEAR}所记1158年的“南宋以宋高宗在位第32年作为诏令、赋役与外交文书的纪年框架”，属于南宋的历史语境。条目所示前因是新岁颁历、年号和纪年是中枢与地方行政沟通的共同前提，其可见后果为为同年政令、战事和地方材料提供日期锚点，不能单独推知具体政策执行。这一变化若进入区域生活，必须通过道路、文书、军队、市场、家庭和地方中介才会发生；政权易手或法令发布不等于所有城镇、乡村与边地在同一时间获得同样经验。

本条所列来源为《宋史》卷24—47〈本纪〉；《建炎以来系年要录》本年条；《宋会要辑稿》相关门类。它们能够支持特定的年序、制度语言、政治行动或局部遗存，却不能无条件化为社会总量。账本特别提醒“纪年框架可靠；该记录仅确证政权纪年连续，年内具体事项须由本年条另证”。因此，军报的兵数、条约的名分、官府的族群分类、判牍中的道德判断以及后出的叙述都须回到产生它们的场景，而不可直接代表全域动机或人口。

将事件与地方层次连接时，应同时考察资源怎样被征集、信息怎样传递、迁徙怎样改变户籍和劳动、宗教与亲属网络怎样提供或限制协作。现有证据往往只照亮少数地点；保留这种不均衡，才能避免把宋辽夏金元的多政权历史改写成单一中心的必然过程。

账本条目\texttt{SS-1159-REGNAL-YEAR}所记1159年的“南宋以宋高宗在位第33年作为诏令、赋役与外交文书的纪年框架”，属于南宋的历史语境。条目所示前因是新岁颁历、年号和纪年是中枢与地方行政沟通的共同前提，其可见后果为为同年政令、战事和地方材料提供日期锚点，不能单独推知具体政策执行。这一变化若进入区域生活，必须通过道路、文书、军队、市场、家庭和地方中介才会发生；政权易手或法令发布不等于所有城镇、乡村与边地在同一时间获得同样经验。

本条所列来源为《宋史》卷24—47〈本纪〉；《建炎以来系年要录》本年条；《宋会要辑稿》相关门类。它们能够支持特定的年序、制度语言、政治行动或局部遗存，却不能无条件化为社会总量。账本特别提醒“纪年框架可靠；该记录仅确证政权纪年连续，年内具体事项须由本年条另证”。因此，军报的兵数、条约的名分、官府的族群分类、判牍中的道德判断以及后出的叙述都须回到产生它们的场景，而不可直接代表全域动机或人口。

将事件与地方层次连接时，应同时考察资源怎样被征集、信息怎样传递、迁徙怎样改变户籍和劳动、宗教与亲属网络怎样提供或限制协作。现有证据往往只照亮少数地点；保留这种不均衡，才能避免把宋辽夏金元的多政权历史改写成单一中心的必然过程。
